% REARIT-ID: REARIT-P001
% REARIT-VERSION: 0.1.0
% REARIT-STATUS: PROPOSTO
% REARIT-TITLE: Princípio da Manutenção Reflexiva Automatizável
% REARIT-DATE: 2026-08-28
% REARIT-ORIGIN: SisTer Infra / OPS-07A0
% REARIT-SUPERSEDES:
% REARIT-SUPERSEDED-BY:

\chapter{REARIT-P001 --- Princípio da Manutenção Reflexiva Automatizável}

\begin{center}
\textbf{Versão 0.1.0 --- PROPOSTO}\\
Origem reflexiva: SisTer Infra / OPS-07A0
\end{center}

\section{Natureza desta formulação}

Este princípio é registrado como formulação emergente da construção do control
plane operacional do SisTer. Ele não deve ser interpretado como afirmação de
que REA ou RIT sempre possuíram, desde suas formulações anteriores, exatamente
este enunciado.

A genealogia relevante é:

\begin{quote}
práticas operacionais inicialmente concretas
$\rightarrow$
reconciliação declarativa
$\rightarrow$
separação entre DEV, LAB e PROD
$\rightarrow$
auditoria do bootstrap histórico
$\rightarrow$
identificação de conhecimento duplicado e manutenção manual
$\rightarrow$
formulação do princípio.
\end{quote}

Sua legitimidade depende tanto de coerência conceitual com REA/RIT quanto de
sua capacidade de explicar e orientar transformações observáveis no SisTer.

\section{Enunciado normativo}

\begin{quote}
\textbf{Toda operação recorrente deve ser automatizável; toda automação deve ser
idempotente quando aplicável; toda transformação deve possuir preflight,
autoridade explícita quando necessária, verificação posterior e evidência; todo
conhecimento que puder ser derivado de declarações autoritativas não deve ser
mantido manualmente em uma segunda representação.}

\medskip

\textbf{A manutenção deve observar invariantes e estado factual, detectar
deriva, inconsistência, obsolescência ou perda de coerência, determinar a menor
ação admissível, executá-la apenas dentro da autoridade disponível, verificar
seus efeitos e preservar a genealogia causal da transição. Quando evidência ou
autoridade forem insuficientes, o sistema deve falhar fechado.}
\end{quote}

\section{Expressão operacional}

\begin{verbatim}
automação
!=
execução cega

automação reflexiva
=
observar
+ comparar
+ explicar
+ validar autoridade
+ agir minimamente
+ verificar
+ registrar evidência
\end{verbatim}

A automação reflexiva não elimina julgamento ou autoridade. Ela torna explícita
a fronteira entre o que pode ser derivado automaticamente e o que exige decisão
legítima externa.

\section{Relação com REA}

A REA organiza a atitude reflexiva em torno da capacidade de perceber,
interpretar e revisar. Aplicada à manutenção, essa orientação implica:

\begin{enumerate}
  \item \textbf{perceber}: observar estado factual e invariantes, sem confundir
  declaração com realidade;
  \item \textbf{interpretar}: comparar o estado observado com a intenção
  autoritativa e explicar a diferença;
  \item \textbf{revisar}: selecionar a menor transformação admissível, dentro da
  autoridade disponível, e reavaliar o estado após agir.
\end{enumerate}

Assim, manutenção não é uma sequência de comandos predefinidos, mas um ciclo
reflexivo governado por evidência.

\section{Relação com RIT}

A Integridade Temporal exige que uma transição preserve simultaneamente:

\begin{itemize}
  \item \textbf{identidade retrospectiva}: ser capaz de reconstruir de onde o
  sistema veio;
  \item \textbf{plasticidade prospectiva}: permitir mudança sem destruir a
  capacidade de evolução;
  \item \textbf{proveniência}: registrar causas, artefatos, decisões e
  evidências relevantes;
  \item \textbf{legitimidade da transição}: distinguir transformação
  tecnicamente possível de transformação efetivamente autorizada.
\end{itemize}

Uma automação que modifica estado sem preservar essas propriedades pode ser
eficiente, mas não é temporalmente íntegra.

\section{Invariantes decorrentes}

\subsection*{P001-I1 --- Automatizabilidade do recorrente}
Operações recorrentes não devem depender de memória procedural humana quando
seus passos puderem ser formalizados e verificados.

\subsection*{P001-I2 --- Idempotência quando aplicável}
A repetição da mesma intenção sobre um estado já convergido não deve produzir
nova transformação material desnecessária.

\subsection*{P001-I3 --- Preflight}
Toda transformação deve validar previamente as condições necessárias para agir.

\subsection*{P001-I4 --- Autoridade explícita}
Transformações que excedam a autoridade do contexto operacional devem falhar
fechado até que autoridade suficiente seja apresentada.

\subsection*{P001-I5 --- Ação mínima}
Entre ações capazes de restabelecer coerência, deve ser escolhida a menor ação
compatível com os invariantes e com a autoridade disponível.

\subsection*{P001-I6 --- Verificação posterior}
Execução bem-sucedida de um comando não constitui evidência suficiente de que a
intenção foi materializada.

\subsection*{P001-I7 --- Evidência}
A transformação deve produzir evidência suficiente para explicar o que foi
observado, decidido, executado e verificado.

\subsection*{P001-I8 --- Não duplicação de autoridade}
Conhecimento derivável de uma fonte declarativa autoritativa não deve ser
mantido manualmente em uma representação concorrente.

\subsection*{P001-I9 --- Fail-closed}
Ausência de evidência ou autoridade necessária deve impedir a transformação,
não ser silenciosamente interpretada como autorização.

\section{Consequências de engenharia}

O princípio favorece arquiteturas baseadas em estado desejado e reconciliação:

\begin{verbatim}
OBSERVE
  |
  v
COMPARE
  |
  v
EXPLAIN
  |
  v
AUTHORITY
  |
  v
MINIMAL ACT
  |
  v
VERIFY
  |
  v
EVIDENCE
\end{verbatim}

Consequências práticas:

\begin{itemize}
  \item check e plan devem ser observacionais;
  \item bootstrap deve materializar pré-condições de forma idempotente;
  \item repair deve atuar apenas sobre deriva comprovada;
  \item apply deve possuir autoridade delimitada pelo modo operacional;
  \item estado já convergido deve resultar em \texttt{NO\_OP};
  \item artefatos derivados devem ser regenerados, não editados como segunda
  fonte de verdade;
  \item transições relevantes devem preservar genealogia causal suficiente para
  auditoria e rollback.
\end{itemize}

\section{Aplicação inicial: SisTer Infra / OPS-07A0}

A auditoria do OPS-07A0 encontrou um bootstrap histórico
(\texttt{setup\_sister\_infra.sh}) que carregava conhecimento concreto de
componentes, portas, hosts e materializações já substituídas por contratos e
declarações autoritativas.

O achado sugere a seguinte fronteira operacional:

\begin{verbatim}
check
  = observar sem corrigir

bootstrap
  = materializar precondicoes locais derivaveis
    de forma idempotente

doctor
  = diagnosticar o ambiente operacional

repair
  = corrigir drift comprovado e autorizado
\end{verbatim}

O bootstrap histórico não deve ser modernizado como segunda implementação do
control plane. Sua história permanece preservada pelo Git, enquanto a
manutenção ativa deve convergir para mecanismos derivados das declarações
canônicas.

\section{Critérios de teste do princípio}

Uma aplicação que reivindique conformidade com REARIT-P001 deve ser capaz de
responder, por evidência:

\begin{enumerate}
  \item Qual fonte define a intenção autoritativa?
  \item Qual estado factual foi observado?
  \item Qual diferença foi detectada?
  \item Por que a ação escolhida é a menor admissível?
  \item Qual autoridade permitiu a transformação?
  \item O que foi efetivamente modificado?
  \item Como o estado posterior foi verificado?
  \item Qual evidência preserva a genealogia da transição?
  \item O que acontece quando evidência ou autoridade são insuficientes?
\end{enumerate}

\section{Hipóteses abertas}

Esta versão ainda não fixa:

\begin{itemize}
  \item um schema universal de evidência operacional;
  \item a granularidade mínima de genealogia exigida para todos os subsistemas;
  \item a relação normativa definitiva entre \texttt{check}, \texttt{doctor},
  \texttt{bootstrap} e \texttt{repair};
  \item quais classes de operação são inerentemente não idempotentes.
\end{itemize}

Esses pontos permanecem explícitos como questões abertas, e não devem ser
preenchidos retrospectivamente como se já estivessem resolvidos.

\section{Referências conceituais}

A formulação dialoga com tradições de cibernética, controle por realimentação,
engenharia de confiabilidade, configuração declarativa, idempotência e
auditabilidade. Entre referências de base para aprofundamento posterior estão
Ashby (\emph{An Introduction to Cybernetics}, 1956), Wiener
(\emph{Cybernetics}, 1948), a disciplina de Site Reliability Engineering e a
semântica de idempotência consolidada em protocolos de sistemas distribuídos.

Essas referências contextualizam o princípio, mas não substituem a genealogia
empírica de sua emergência dentro do SisTer.
